\documentclass[11pt,a4paper]{article}

\usepackage[T1]{fontenc}
\usepackage[margin=1in]{geometry}
\usepackage{xcolor}
\usepackage{listings}

\definecolor{codegreen}{rgb}{0,0.45,0}
\definecolor{codegray}{rgb}{0.4,0.4,0.4}
\definecolor{codepurple}{rgb}{0.5,0,0.5}

\lstdefinestyle{cpp}{
    language=C++,
    basicstyle=\ttfamily\small,
    keywordstyle=\color{blue},
    commentstyle=\color{codegreen},
    stringstyle=\color{codepurple},
    numbers=left,
    numberstyle=\tiny\color{codegray},
    stepnumber=1,
    numbersep=8pt,
    breaklines=true,
    breakatwhitespace=false,
    keepspaces=true,
    showstringspaces=false,
    tabsize=4,
    frame=single
}

\title{B-Tree Implementation in C++}
\author{Advanced Data Structures}
\date{}

\begin{document}

\maketitle

\section{Overview}

This document collects the B-Tree implementation and its example program.
The C++ source files remain separate so the project can still be built and
maintained normally. Compile this document from the repository root.

\section{Node Types}

\subsection{\texttt{include/Node.h}}
\lstinputlisting[style=cpp,caption={Base node}]{include/Node.h}

\subsection{\texttt{include/LeafNode.h}}
\lstinputlisting[style=cpp,caption={Leaf node}]{include/LeafNode.h}

\subsection{\texttt{include/InternalNode.h}}
\lstinputlisting[style=cpp,caption={Internal node}]{include/InternalNode.h}

\section{B-Tree}

\subsection{\texttt{include/BTree.h}}
\lstinputlisting[style=cpp,caption={B-Tree declaration}]{include/BTree.h}

\subsection{\texttt{include/BTree.tpp}}
\lstinputlisting[style=cpp,caption={B-Tree template implementation}]{include/BTree.tpp}

\section{Example Program}

\subsection{\texttt{src/main.cpp}}
\lstinputlisting[style=cpp,caption={Example program and basic operation checks}]{src/main.cpp}

\end{document}
